\section{Análisis temporal}
\label{sec:latencia}

La tabla~\ref{tab:latencia} estima el retardo desde que el operador mueve la mano hasta que el servo recibe la nueva orden. \textbf{Son estimaciones a partir del código y de las hojas de datos, no mediciones} sobre el hardware.

\begin{table}[H]
\centering
\begin{tabularx}{\textwidth}{X c X}
\toprule
\textbf{Etapa} & \textbf{Retardo} & \textbf{Origen} \\
\midrule
Espera hasta el siguiente ciclo del transmisor & 0--\SI{20}{\milli\second} (media 10) & Muestreo a \SI{50}{\hertz} \\
Filtro interno del MPU6050 (sólo muñeca) & $\approx\SI{5}{\milli\second}$ & DLPF de \SI{44}{\hertz} \cite{mpu6050_regmap} \\
Lectura de dedos, cálculo y empaquetado & $\approx\SI{5}{\milli\second}$ & Sobremuestreo $16\times$ \\
Transmisión ESP-NOW & $\approx$ 1--\SI{6}{\milli\second} & Valores típicos reportados \cite{espnow_faq, espnow_issue115} \\
Tarea de control del receptor & $<\SI{1}{\milli\second}$ & Despierta al llegar el paquete \\
Siguiente pulso PWM del servo & 0--\SI{20}{\milli\second} (media 10) & PWM de \SI{50}{\hertz} \\
Limitador de velocidad del receptor & 0--$\approx$\SI{160}{\milli\second} & \SI{8}{\degree} por ciclo en los dedos (tabla~\ref{tab:param_rx}): 0 para cambios chicos, unos 8 ciclos para un recorrido completo \\
Movimiento mecánico & según amplitud & $\approx\SI{0.1}{\second}$ cada \SI{60}{\degree} (SG90 sin carga) \\
\bottomrule
\end{tabularx}
\caption{Estimación de la cadena de retardos (versión 2.2.0 y posteriores).}
\label{tab:latencia}
\end{table}

En promedio, la orden llega al servo unos \SI{30}{\milli\second} después del movimiento del operador. Ante cambios grandes se suma el limitador de velocidad, que reparte el movimiento en varios ciclos, y luego el tiempo que tarda el servo en recorrer el ángulo.

Para medir la latencia real se preparó un ensayo de punta a punta sobre el camino de los dedos (procedimiento en \texttt{docs/\allowbreak pruebas/\allowbreak PRUEBA\_LATENCIA.md}). El receptor reemplaza al potenciómetro del índice: aplica un escalón de tensión en la entrada del transmisor (GPIO 34) y mide, con su propio reloj, el tiempo hasta el fin del primer pulso PWM con el ancho nuevo en la salida del servo del índice. Como el inicio y el fin se toman en la misma placa, no hace falta sincronizar relojes. La medición incluye la espera al ciclo del transmisor, la lectura y el envío, la recepción, la tarea de control y la espera al pulso PWM; no incluye el movimiento mecánico ni el filtro interno del MPU6050.

\pendiente{Faltan los resultados de este ensayo: con ellos se confirmará o corregirá esta estimación.}
